% assumption.tex — 二、模型假设（2026-08-15, 每条编号并附依据）

为建立可求解的数学模型，结合题目条件与数据特征作出如下假设：

\begin{enumerate}[label=（\arabic*）]
\item 附件中的销售流水、批发价格与损耗率数据真实可靠，无系统性缺失。
  依据：数据覆盖三个完整年度且日均交易约 810 笔，量级与商超实际经营相符；
  预处理中发现的退货记录（461 条，占 0.05\%）予以剔除，其余记录保留。

\item 未来一周的销量与批发价由星期效应主导，可用历史 6--7 月同星期均值预测。
  依据：对周内销量的统计显示周六周日销量明显高于工作日，品类日销量序列存在以周
  为周期的规律；同期（6--7 月）数据与预测时段季节一致。

\item 各蔬菜单品独立销售，不存在捆绑促销与跨品类替代。
  依据：附件销售流水按单品逐笔记录，无组合售卖字段；问题一中计算的关联关系仅用于
  分析，不进入补货模型。

\item 损耗率在决策周内保持近期水平；当日未售出部分按历史折扣水平打折清仓。
  依据：附件4 的损耗率由近期盘点数据计算；附件2 中打折交易占比稳定，可用各品类
  历史打折均价与正常均价的比值刻画折扣水平。

\item 商超日展示容量可用历史 6--7 月日总销量的 95\% 分位估计（下限 67.5 kg）。
  依据：某日实现的销售量意味着当日货架至少承载了同等规模的商品；取 95\% 分位可覆盖
  绝大多数营业日的需求水平，67.5 kg 为 27 个单品与最小陈列量 2.5 kg 的乘积，保证
  约束可行。

\item 各品类销量与成本加成定价的关系可用三次多项式在观测价格区间内近似。
  依据：对线性、二次、三次等候选形式做拟合比较，三次形式的拟合优度整体最高
  （$R^2$ 最高 0.58），且能刻画低价区间需求刚性、高价区间需求下降的非对称形态。
\end{enumerate}
